\chapter{Cronograma}
\label{chap:cronograma}

O cronograma de execução foi organizado para o próximo semestre letivo, iniciando em março de 2026 e finalizando em julho de 2026, contemplando as etapas necessárias para a avaliação de algoritmos de agregação robusta em Aprendizado Federado sob cenários adversos. [file:40]
A revisão bibliográfica e a consolidação dos conceitos fundamentais (Aprendizado Federado, ataques e agregação robusta) já foram realizadas, restando as etapas de implementação, experimentação, análise e escrita final do trabalho. [file:40]

\begin{table}[htbp]
    \centering
    \caption{Cronograma de execução (março a julho de 2026).}
    \label{tab:cronograma}
    \begin{tabular}{|p{8.2cm}|c|c|c|c|c|}
        \hline
        \textbf{Etapa} & \textbf{Mar./2026} & \textbf{Abr./2026} & \textbf{Mai./2026} & \textbf{Jun./2026} & \textbf{Jul./2026} \\
        \hline
        Planejamento experimental (cenários, seeds, métricas e protocolo) & X &  &  &  &  \\
        \hline
        Preparação de dados e particionamento federado (IID e Non-IID) & X & X &  &  &  \\
        \hline
        Implementação do ambiente federado e do baseline (FedAvg) & X & X &  &  &  \\
        \hline
        Implementação dos agregadores robustos (ex.: Krum, Median, Trimmed Mean e/ou modernos) &  & X & X &  &  \\
        \hline
        Implementação dos ataques (ex.: sign-flip, scaling, MinMax e adaptativos) &  & X & X &  &  \\
        \hline
        Execução dos experimentos (variação de heterogeneidade e fração de atacantes) &  &  & X & X &  \\
        \hline
        Avaliação e análise dos resultados (curvas, tabelas, robustez e overhead) &  &  &  & X & X \\
        \hline
        Escrita do TCC (metodologia, resultados e discussão) &  &  &  & X & X \\
        \hline
        Revisão final, normalização ABNT e ajustes de texto/figuras &  &  &  &  & X \\
        \hline
        Preparação da apresentação (slides e ensaio) e entrega/defesa &  &  &  &  & X \\
        \hline
    \end{tabular}

    \vspace{0.3cm}
    {\footnotesize Fonte: Elaborado pelo autor.}
\end{table}

As atividades foram distribuídas de forma a priorizar a implementação incremental do ambiente federado, iniciando pelo baseline e, em seguida, incorporando agregadores robustos e ataques adversariais. [file:40]
Após a etapa de experimentação, os resultados serão analisados de acordo com as métricas definidas no protocolo experimental descrito no trabalho, permitindo comparar desempenho e robustez sob diferentes níveis de heterogeneidade e adversarialidade. [file:40]




% \chapter{Cronograma}
% \label{chap:cronograma}
%
% O cronograma de execução do projeto foi estruturado de modo a organizar as atividades ao longo de dois meses, abrangendo as principais etapas necessárias para o \gls{TCC}. Além disso, a revisão bibliográfica do tema e dos conceitos abordados neste trabalho já foram realizados.  As atividades programadas foram distribuídas entre janeiro e fevereiro de 2025. 
%
% \begin{table}[h]
%     \centering
%     \caption{Planejamento das etapas seguintes.}
%     \label{tab:planejamento}
%     
%     \begin{longtable}{|l|c|c|}
%     \hline
%     \textbf{Etapa} & \textbf{Janeiro 2025} & \textbf{Fevereiro 2025} \\
%     \hline
%     \endfirsthead
%     \hline
%     \textbf{Etapa} & \textbf{Janeiro 2025} & \textbf{Fevereiro 2025} \\
%     \hline
%     \endhead
%     \hline
%     Coleta e Processamento de Dados & X & - \\
%     Implementação dos Modelos & - & X \\
%     Avaliação e Análise & - & X \\
%     Escrita do TCC & - & X \\
%     Revisão do Texto & - & X \\
%     Apresentação do TCC & - & X \\
%     \hline
%     \end{longtable}
%
%     \parbox{\linewidth}{\raggedright \footnotesize Fonte: Próprio autor.}
% \end{table}
%
% A etapa de Coleta e Processamento de Dados tem como objetivo reunir todos os dados necessários para a análise e preparar as informações para as fases seguintes do projeto. Além disso, será realizada a Implementação dos Modelos, que são essenciais para a análise dos dados coletados. Durante esse período, também será feita os ajustes necessários para a execução correta desses algoritmos. Após a implementação dos modelos, será iniciada a Avaliação e Análise dos Resultados, verificando e analisando dos dados processados e comparando os resultados obtidos pelos algoritmos.  A Escrita do \gls{TCC} como o próprio nome sugere inclui a elaboração dos capítulos do trabalho, como introdução, metodologia, resultados e discussão. Após a redação do \gls{TCC}, será realizada uma Revisão do Texto, focando em aspectos como clareza, coesão, ortografia e formatação. Por fim, será realizada a Apresentação do \gls{TCC} que incluirá a organização da defesa, elaboração de slides e a exposição oral do trabalho. O cronograma apresentado na Tabela 1 demonstra as etapas que serão realizadas durante a elaboração do \gls{TCC}.